% TOP_PROBLEM: {"id": 30006848, "title": "The Dynamic Optimality Conjecture for Splay Trees", "category_id": 15, "category": {"id": 15, "name": "computer_science", "display_name": "Computer Science", "description": "Computational complexity, algorithms, and theoretical CS.", "slug": "computer-science", "order_index": 15, "created_at": "2026-07-31T15:26:25.671Z"}, "status": "open"}
% FIRST_POSTED: 2026-09-27
% !TeX program = lualatex
% ATTEMPT_STATUS: unresolved
\documentclass[11pt]{article}
\usepackage[margin=25mm]{geometry}
\usepackage{fontspec}
\setmainfont{Latin Modern Roman}
\usepackage{amsmath,amssymb,mathtools}
\usepackage{unicode-math}
\setmathfont{Latin Modern Math}
\usepackage{xurl,hyperref}
\hypersetup{colorlinks=true,urlcolor=blue}
\setcounter{secnumdepth}{0}
\setlength{\parindent}{0pt}
\setlength{\parskip}{5pt}
\setlength{\emergencystretch}{3em}
\begin{document}
\section*{\texorpdfstring{The Dynamic Optimality Conjecture for Splay Trees}{The Dynamic Optimality Conjecture for Splay Trees}}\label{30006848}
\subsection{Definitions and mathematical statement}
In the dynamic binary-search-tree model, a fixed ordered set of $n$ keys is stored in a rooted binary search tree. Servicing an access permits pointer moves and rotations, each charged unit cost under the source's conventions. Splaying moves the accessed key to the root by the standard zig, zig--zig, and zig--zag steps. For an access sequence $\sigma$, let $\operatorname{OPT}(\sigma)$ be the minimum cost over offline BST algorithms, which may know the complete sequence in advance. Dynamic optimality asks for an absolute constant $C$ such that
\[
 \operatorname{cost}_{\rm splay}(\sigma)
 \le C\operatorname{OPT}(\sigma)+I(n)
\]
for every $\sigma$, where $I(n)$ is the standard initialization allowance in the chosen model, independent of the access sequence. The input catalogue does not fix a numeric expression for $I$; a precise comparison must retain the cited source's initial-tree convention, rather than allowing an arbitrary sequence-dependent additive error.
\par\smallskip\textit{Formulation and concept references:} \href{https://arxiv.org/abs/2607.18498}{[S1]}.

\subsection{Short English statement}
Are splay trees always within a universal constant factor of the best binary-search-tree strategy that already knows the entire access sequence?
\subsection{Research attempt}
\textbf{Status: UNRESOLVED.} A complete weighted access-lemma argument is given, yielding logarithmic competitiveness and a constant-factor comparison with a fixed search tree, each with an explicit initialization term. The argument also handles repeated access to one key optimally. It does not compare splaying within a constant factor to every changing offline tree.

\paragraph{Source audit and initialization convention.}
The July 2026 manuscript [S1], consulted on 27 September 2026, claims an $O(\log\log n\cdot\log^2\log\log n)$ competitive ratio. That factor still grows with $n$ and is not the conjectured absolute constant. Its model section charges one plus the number of rotations, requires the accessed key to end at the root, allows the offline tree an optimal initial configuration, and treats the splay tree's initial configuration as arbitrary. Its appendix explains equivalence with the usual BST model. The full new competitive analysis is not independently reproved here.

The argument below works in that root-and-rotation cost convention. Search along the path and ordinary splaying in the pointer model cost a comparable number of operations, but precise initialization terms must still be retained. In particular, an additive term depending on the length or content of the access sequence is not a permissible substitute for the catalogue's fixed $I(n)$.

\paragraph{1. Weights, ranks and the potential.}
Assign each key a fixed positive weight $w_x$. For a current tree, let $s(x)$ be the sum of weights in the subtree rooted at $x$, let $r(x)=\log_2s(x)$, and set
\[
 W=\sum_xw_x,\qquad \Phi=\sum_x r(x).
\]
All weights stay fixed during the analysis. A rotation changes subtree sets only for the rotated nodes; the subtree of the root of the affected local configuration has the same total weight before and after the operation. This local conservation permits an exact bound on the potential change.

Let $x$ be the accessed key, $y$ its parent and, where present, $z$ its grandparent. Primes denote quantities after the current splay step. Each double step costs two rotations, and a final zig costs one.

\paragraph{2. The single zig.}
When $y$ is the root of the affected tree and $x$ is rotated above it, $s'(x)=s(y)$. Only the ranks of $x,y$ change, so
\[
 1+\Delta\Phi=1+r'(y)-r(x)
 \le1+r'(x)-r(x).
\]
The same argument applies to either left or right orientation. This step occurs at most once in a splay, at the end.

\paragraph{3. The zig-zig step.}
Suppose $x$ and $y$ are both left children; the reflected case is identical. Rotate $y$ above $z$ and then $x$ above $y$. Since $r'(x)=r(z)$,
\[
 2+\Delta\Phi=2+r'(y)+r'(z)-r(x)-r(y).
\]
We have $r'(y)\le r'(x)$ and $r(y)\ge r(x)$. Moreover the old subtree of $x$ and the new subtree of $z$ are disjoint subsets of the new subtree of $x$. Thus
\[
 s(x)+s'(z)\le s'(x),\qquad
 r(x)+r'(z)\le2r'(x)-2,
\]
where the latter inequality follows from $ab\le((a+b)/2)^2$. Substitution yields
\[
 2+\Delta\Phi\le3\bigl(r'(x)-r(x)\bigr).
\]
The disjointness statement concerns these particular old and new subtrees; using arbitrary same-named subtrees on the two sides would not justify the estimate.

\paragraph{4. The zig-zag step.}
Suppose $x$ is the right child of a left child $y$ of $z$. Rotate $x$ above $y$ and then above $z$. Again $r'(x)=r(z)$ and the same three-rank expression gives the potential change. Now the new subtrees of $y,z$ are disjoint subsets of the new subtree of $x$, so
\[
 r'(y)+r'(z)\le2r'(x)-2.
\]
Together with $r(y)\ge r(x)$ this gives
\[
 2+\Delta\Phi\le2\bigl(r'(x)-r(x)\bigr)
 \le3\bigl(r'(x)-r(x)\bigr).
\]
The rank difference is nonnegative because the accessed node's subtree grows. The other zig-zag orientation follows by reflection.

Summing all the steps telescopes the ranks of $x$. If $R$ rotations are performed and $s_0(x)$ was the initial subtree weight of the accessed key, then
\[
 R+\Phi_{\rm after}-\Phi_{\rm before}
 \le3\log_2\frac W{s_0(x)}+1
 \le3\log_2\frac W{w_x}+1.
\]
The root-and-rotation access cost is $R+1$, so its amortized cost is at most $3\log_2(W/w_x)+2$. This also holds for an access already at the root.

\paragraph{5. Uniform weights give a weaker competitive theorem.}
Choose $w_x=1$. Every subtree size lies between one and $n$, so $0\le\Phi\le n\log_2n$. For a sequence of $m$ accesses,
\[
 \operatorname{cost}_{\rm splay}
 \le m(3\log_2n+2)+n\log_2n.
\]
Every competing strategy pays at least one per access, hence $\operatorname{OPT}\ge m$. This proves an $O(\log n)$ competitive factor with an explicit $O(n\log n)$ initialization allowance. It is weaker both than the claimed recent result and than the desired constant factor.

The displayed additive term is a bound supplied by this potential proof; it does not redefine the unspecified $I(n)$ in the catalogue or assert that it meets a sharper initialization convention in another theorem. Its independence from $m$ is what makes it a legitimate initialization term for this weaker comparison.

\paragraph{6. A constant-factor comparison with a fixed reference tree.}
Let $T_*$ be any fixed BST, and let $d_x$ be the depth of key $x$ in it. Give $x$ weight $w_x=4^{-d_x}$. Since a binary tree has at most $2^d$ nodes at depth $d$,
\[
 W\le\sum_{d\ge0}2^d4^{-d}=2,
 \qquad
 \log_2(W/w_x)\le1+2d_x.
\]
For these fixed weights, the possible range of $\Phi$ is at most
\[
 \sum_x\log_2(W/w_x)\le n+2\sum_xd_x\le2n^2.
\]
Therefore
\[
 \operatorname{cost}_{\rm splay}(\sigma)
 \le6\sum_{x\text{ accessed in }\sigma}(d_x+1)+2n^2.
\]
The sum is the ordinary search cost in the fixed reference tree. The access frequencies need not be known to splaying; they do not enter the algorithm. This establishes a static comparison with a sequence-independent additive bound even if the reference tree is selected after observing the sequence.

It does not establish dynamic optimality. The best fixed reference tree may cost far more than the best changing tree. Changing the weights after each access to track an offline comparator would create additional potential changes, and those changes have not been charged by the access lemma.

\paragraph{7. An explicit gap between static and dynamic comparisons.}
Consider many repetitions of the ordered sweep $1,2,\ldots,n$. In any fixed binary search tree, the sum of all key depths is $\Omega(n\log n)$: at depth less than $h$ there are at most $2^h-1$ nodes, and summing the resulting depth lower bounds gives this estimate. Thus the fixed-tree search cost per full sweep is $\Omega(n\log n)$.

A changing tree can serve these sweeps in $O(n)$ rotations per sweep. Start with a right spine rooted at 1. After accessing $i$, rotate its right child $i+1$ to the root for the next access. At the end the tree is a left spine rooted at $n$. Before the next access to 1, rotate successive left children to the root, restoring the right spine with $n-1$ rotations. Together with the $n$ unit access charges, each complete sweep costs at most $3n-2$ after the chosen initial setup. The offline comparator is allowed that initial shape in the source convention.

This is not a bad sequence for splaying and is not a counterexample to dynamic optimality. It shows why a theorem comparing to the best fixed tree cannot simply be relabelled as a theorem comparing to the best changing tree: the two comparators can differ by a logarithmic factor.

For repeated access to one key $x$, a direct argument does give a constant comparison. The first splay uses at most $n-1$ rotations; every later access finds $x$ at the root and costs one. Thus the whole sequence costs at most $m+n-1$, whereas every competitor pays at least $m$. This fully handles that restricted access family.

\paragraph{8. Exact finite diagnostics and the missing amortization.}
The offline script enumerates all BST shapes through five keys and every accessed key, with uniform and specified nonuniform positive integer weights. It performs the actual zig, zig-zig and zig-zag transformations and checks the access inequality after exponentiating it to an exact rational product inequality. It also checks the repeated-key and ordered-sweep constructions. No floating logarithm is used as a proof certificate for the access lemma.

For an additional bounded comparison it enumerates all access sequences of length at most four on at most four keys. A shortest-path computation on the finite rotation graph supplies the exact offline optimum in the source's root-and-rotation model, including free choice of its initial tree. These finite optima are diagnostics only; neither their ratios nor the weighted access lemma imply an absolute competitive factor for all $n$ and all sequence lengths.

The unresolved step is an amortized comparison that charges every kind of splay restructuring to an unrestricted changing optimum with only a universal constant loss and the prescribed initialization convention. The fixed-weight potential above supplies no such charge for changes in the reference tree, so the catalogue conjecture remains unresolved here.

\subsection{Sources}
\begin{itemize}
\item[S1]Petr Chmel et al., Splay trees are almost dynamically optimal (2026).
\url{https://arxiv.org/abs/2607.18498}.
\textit{Formulation source.}
\end{itemize}
\end{document}
